\documentclass[11pt]{article}
% ===========================================================================
%  Shared preamble for the AP Calculus AB exam-preparation worksheets.
%  Input right after \documentclass. Files beginning with "_" are not compiled
%  on their own. Navy and gold, to match the course's lesson pages.
% ===========================================================================
\usepackage[letterpaper,top=2.3cm,bottom=2.2cm,left=1.7cm,right=1.7cm,headheight=28pt,footskip=30pt]{geometry}
\usepackage{amsmath,amssymb}
\usepackage[T1]{fontenc}
\usepackage[utf8]{inputenc}
\usepackage{lmodern}
\usepackage{xcolor}
\usepackage[most]{tcolorbox}
\usepackage{enumitem}
\usepackage{fancyhdr}
\usepackage{array}
\usepackage{tikz}
\usepackage{pgfplots}
\pgfplotsset{compat=1.18}

\definecolor{navy}{HTML}{161E54}
\definecolor{ink}{HTML}{27357E}
\definecolor{gold}{HTML}{B98A00}
\definecolor{goldbg}{HTML}{FFFDF3}
\definecolor{inkbg}{HTML}{F6F7FD}
\definecolor{rule}{HTML}{C9CFF0}
\definecolor{okgreen}{HTML}{1B7A44}
\definecolor{okbg}{HTML}{F1F8F3}

\pagestyle{fancy}
\fancyhf{}
\fancyhead[L]{\footnotesize Name: \underline{\hspace{4cm}}}
\fancyhead[C]{\textbf{\color{navy}Integration Academy}\\[-2pt]{\scriptsize Calculus AB \textperiodcentered{} AP\textsuperscript{\textregistered} Exam Preparation}}
\fancyfoot[L]{\tiny\color{black!55}\parbox[t]{0.78\textwidth}{AP\textsuperscript{\textregistered} is a registered trademark of the College Board, which was not involved in the production of, and does not endorse, this material. \textcopyright{} Integration Academy.}}
\fancyfoot[R]{\scriptsize\thepage}
\renewcommand{\headrulewidth}{1.2pt}
\renewcommand{\headrule}{\hbox to\headwidth{\color{gold}\leaders\hrule height \headrulewidth\hfill}}
\renewcommand{\footrulewidth}{0pt}
\setlength{\parindent}{0pt}
\setlength{\parskip}{4pt}

% Title block: \wstitle{code}{title}{unit line}
\newcommand{\wstitle}[3]{\begin{center}{\small\bfseries\color{gold}WORKSHEET #1}\\[3pt]{\LARGE\bfseries\color{navy}#2}\\[3pt]{\small\color{black!60}#3}\end{center}\vspace{2pt}}
% Section heading
\newcommand{\sech}[1]{\par\vspace{10pt}{\large\bfseries\color{navy}#1}\par\nobreak\vspace{2pt}{\color{rule}\hrule height0.8pt}\vspace{7pt}}
% A level tag printed beside a group of questions
\newcommand{\level}[1]{\par\vspace{6pt}{\footnotesize\bfseries\color{gold}\MakeUppercase{#1}}\par\nobreak\vspace{2pt}}

\newtcolorbox{ideabox}{colframe=ink,colback=inkbg,coltitle=white,colbacktitle=ink,fonttitle=\bfseries\small,title={Key ideas},arc=2pt,boxrule=0.8pt,left=8pt,right=8pt}
\newtcolorbox{example}[1]{colframe=rule,colback=white,coltitle=navy,colbacktitle=inkbg,fonttitle=\bfseries\small,title={#1},arc=2pt,boxrule=0.8pt,borderline west={3pt}{0pt}{ink},left=10pt,right=8pt}
\newtcolorbox{frbox}[1]{colframe=gold,colback=goldbg,coltitle=white,colbacktitle=gold,fonttitle=\bfseries\small,title={#1},arc=2pt,boxrule=0.8pt,breakable,left=8pt,right=8pt}
\newtcolorbox{solbox}{colframe=okgreen,colback=okbg,coltitle=white,colbacktitle=okgreen,fonttitle=\bfseries\small,title={Worked solutions},arc=2pt,boxrule=0.8pt,breakable,left=8pt,right=8pt}
\newtcolorbox{scorebox}[1]{colframe=navy,colback=white,coltitle=white,colbacktitle=navy,fonttitle=\bfseries\small,title={#1},arc=2pt,boxrule=0.8pt,breakable,left=8pt,right=8pt}

\newcommand{\soln}{\par\smallskip\textbf{\color{ink}Solution.}\ }
\newcommand{\result}[1]{\par\smallskip\textbf{#1}}
% blank working space of a given height
\newcommand{\work}[1]{\par\vspace{#1}}
% a point value in the margin of a scoring guide
\newcommand{\pts}[1]{\hfill{\footnotesize\bfseries\color{gold}[#1]}}

% Coordinate grid for graph questions: \begin{apgraph}{xmin}{xmax}{ymin}{ymax} ... \end{apgraph}
\newenvironment{apgraph}[4]{\begin{center}\begin{tikzpicture}\begin{axis}[
  axis lines=middle, xmin=#1, xmax=#2, ymin=#3, ymax=#4,
  grid=both, grid style={black!12}, xtick distance=1, ytick distance=1,
  ticklabel style={font=\scriptsize}, xlabel={$x$}, ylabel={$y$},
  width=10cm, height=7.2cm, samples=80, clip=false,
  every axis plot/.append style={line width=1.1pt, color=ink}]}{\end{axis}\end{tikzpicture}\end{center}}
\newcommand{\closeddot}[2]{\addplot[only marks, mark=*, mark size=2.2pt, color=ink] coordinates {(#1,#2)};}
\newcommand{\opendot}[2]{\addplot[only marks, mark=*, mark size=2.2pt, color=ink, mark options={fill=white}] coordinates {(#1,#2)};}

% Multiple-choice options: four across, or stacked when the options are long.
\newcommand{\opts}[4]{\par\vspace{3pt}\begin{tabular}{@{}p{0.235\linewidth}p{0.235\linewidth}p{0.235\linewidth}p{0.235\linewidth}@{}}(A)\ #1 & (B)\ #2 & (C)\ #3 & (D)\ #4\end{tabular}\par\vspace{9pt}}
\newcommand{\optsv}[4]{\par\vspace{3pt}\begin{tabular}{@{}l@{\ \ }p{0.9\linewidth}@{}}(A) & #1\\(B) & #2\\(C) & #3\\(D) & #4\end{tabular}\par\vspace{9pt}}

\fancyhead[R]{\footnotesize\color{black!60}Worksheet 7.6}
\begin{document}

\wstitle{7.6}{Particular Solutions and Their Domains}{Unit 7: Differential Equations}

\begin{ideabox}
\begin{itemize}[leftmargin=*,itemsep=2pt,topsep=2pt]
  \item Use the initial condition to find the constant \textbf{before} solving for $y$ where possible; the algebra is easier.
  \item When a square root appears, the initial condition decides the sign.
  \item When an absolute value appears, the sign of the quantity at the initial point decides how to remove it.
  \item The \textbf{domain} of a particular solution is the largest single interval, containing the initial $x$-value, on which the solution is defined and satisfies the equation.
\end{itemize}
\end{ideabox}

\sech{Worked examples}
\begin{example}{Example 1 \textperiodcentered{} An exponential}
Solve $\dfrac{dy}{dx}=\dfrac y3$ with $y(0)=6$.\soln $\ln|y|=\dfrac x3+C$, so $y=Ae^{x/3}$ with $A=6$.\result{$y=6e^{x/3}$.}
\end{example}

\begin{example}{Example 2 \textperiodcentered{} Choosing a sign}
Solve $\dfrac{dy}{dx}=\dfrac{4x^3}{y}$ with $y(0)=-2$.\soln $\dfrac{y^2}{2}=x^4+C$ with $C=2$, so $y^2=2x^4+4$. The initial value is negative.\result{$y=-\sqrt{2x^4+4}$.}
\end{example}

\begin{example}{Example 3 \textperiodcentered{} A domain}
Solve $\dfrac{dy}{dx}=\dfrac{y^2}{x^2}$ with $y(1)=\dfrac12$, and state the domain.\soln $-\dfrac1y=-\dfrac1x+C$; at $\left(1,\frac12\right)$: $-2=-1+C$, so $C=-1$.\result{$y=\dfrac{x}{x+1}$, for $x>0$ (the equation is undefined at $x=0$).}
\end{example}

\newpage
\sech{Practice}

\level{Core}
\begin{enumerate}[leftmargin=*,itemsep=2pt]
  \item Solve $\dfrac{dy}{dx}=-4y$ with $y(0)=3$.\work{2.3cm}
  \item Solve $\dfrac{dy}{dx}=2x$ with $y(1)=5$.\work{2.3cm}
  \item Solve $\dfrac{dy}{dx}=\cos x$ with $y(0)=-1$.\work{2.3cm}
  \item Solve $\dfrac{dy}{dx}=\dfrac1x$ for $x>0$ with $y(1)=2$.\work{2.3cm}
\end{enumerate}

\level{Exam level}
\begin{enumerate}[leftmargin=*,itemsep=2pt,start=5]
  \item Solve $\dfrac{dy}{dx}=\dfrac{9x}{y}$ with $y(0)=-3$.\work{2.8cm}
  \item Solve $\dfrac{dy}{dx}=3(y-2)$ with $y(0)=5$.\work{2.8cm}
  \item Solve $\dfrac{dy}{dx}=e^{y}$ with $y(0)=0$, and state the domain of the solution.\work{2.8cm}
  \item Solve $\dfrac{dy}{dx}=y\cos x$ with $y(0)=4$.\work{2.8cm}
\end{enumerate}

\begin{frbox}{Question 9 \textperiodcentered{} Free-response practice}
Consider the differential equation $\dfrac{dy}{dx}=x(y-1)^2$. Let $y=f(x)$ be the particular solution with $f(0)=3$.
\begin{enumerate}[label=(\alph*),leftmargin=*,itemsep=2pt]
  \item Find the slope of the graph of $f$ at $(0,3)$. What does this tell you about the graph there?\work{2.6cm}
  \item Find the particular solution $y=f(x)$.\work{2.6cm}
  \item State the domain of the solution, and give a reason.\work{2.6cm}
\end{enumerate}
\end{frbox}

\newpage
\begin{solbox}
\begin{enumerate}[leftmargin=*,itemsep=4pt]
  \item $y=3e^{-4x}$.
  \item $y=x^2+C$ with $C=4$: $y=x^2+4$.
  \item $y=\sin x-1$.
  \item $y=\ln x+2$.
  \item $y^2=9x^2+C$ with $C=9$; the initial value is negative: $y=-3\sqrt{x^2+1}$.
  \item $\ln|y-2|=3x+C$; $y-2=3$ at the start: $y=2+3e^{3x}$.
  \item $-e^{-y}=x+C$ with $C=-1$: $e^{-y}=1-x$, so $y=-\ln(1-x)$, for $x<1$.
  \item $\ln|y|=\sin x+C$: $y=4e^{\sin x}$.
  \item \begin{enumerate}[label=(\alph*),leftmargin=*,itemsep=3pt]
    \item $0\cdot(2)^2=0$: the tangent is horizontal at $(0,3)$.
    \item $\dfrac{dy}{(y-1)^2}=x\,dx$ gives $-\dfrac{1}{y-1}=\dfrac{x^2}{2}+C$, and $C=-\dfrac12$. So $y-1=\dfrac{2}{1-x^2}$, that is $y=1+\dfrac{2}{1-x^2}$.
    \item $-1<x<1$: the solution is undefined at $x=\pm1$, and the domain is the single interval containing $x=0$.
  \end{enumerate}
\end{enumerate}
\end{solbox}

\end{document}
